\documentclass[runningheads]{llncs}

\usepackage[T1]{fontenc}
\usepackage{graphicx}
\usepackage[cmex10]{amsmath}
\usepackage{amssymb}
\usepackage{booktabs}

\newcommand{\cmark}{\checkmark}
\newcommand{\xmark}{\ensuremath{\times}}

\begin{document}

\title{A Model-Agnostic Temporal Explainability Framework\\
for Financial Time Series Prediction}

\titlerunning{Model-Agnostic Temporal Explainability for Finance}

\author{Rushell Vanessa Zavalaga Orozco\inst{1} \and
Juan Carlos Gutiérrez Cáceres\inst{1} \and
Ana María Cuadros Valdivia\inst{1}}

\authorrunning{R. V. Zavalaga Orozco et al.}

\institute{School of Computer Science,\\
Universidad Nacional de San Agustín de Arequipa, Arequipa, Peru\\
\email{rzavalaga@unsa.edu.pe, jgutierrezca@unsa.edu.pe, acuadrosv@unsa.edu.pe}}

\maketitle

% =======================================================
\begin{abstract}
Stock market forecasting models based on deep learning achieve high
accuracy but function as black boxes: they produce predictions without
revealing which historical observations---which variables, on which
days---drove them. Without knowing which inputs drive a prediction,
it is impossible to validate model behavior, detect spurious
correlations, or determine whether different architectures learn the
same representations from financial data.

Existing explanation methods have critical gaps: some require access
to the model's internal computations, making them inapplicable to
non-neural models; others ignore the sequential structure of financial
data, generating explanations based on combinations of observations
that could never occur in real markets.

This paper proposes a framework that explains any forecasting model
by systematically asking: \textit{``if I had not known the value of
variable $v$ on day $t$, how different would the prediction have
been?''} The answer is an importance matrix
$\hat{\alpha} \in \mathbb{R}^{T \times V}$ showing each variable's
contribution at each historical day. The framework communicates with
the model only through its prediction function, requiring no internal
access to weights, gradients, or architecture.

Applied to three structurally different models---LSTM, Transformer,
and Random Forest---trained on S\&P~500 data (2010--2024): the
framework correctly identifies which inputs matter when the ground
truth is known (AUP\,=\,0.9960), runs at 17.6\,ms per
explanation---up to $160{\times}$ faster than competing
methods---and produces explanations that differ appropriately between
model types (Spearman $\rho=0.748$ between neural models vs.\
$\rho<0.40$ against the ensemble model, $p<10^{-300}$), confirming
that the framework captures model-specific behavior rather than
artifacts of the explanation procedure. This pattern is robust to the
choice of reference vector, to the market regime evaluated, and to a
naturally bounded target formulation that also shows none of the
three models beats a naive persistence baseline, in magnitude or in
direction, consistent with weak-form market efficiency.

\keywords{Financial time series \and Post-hoc explainability \and
Model-agnostic \and Temporal perturbations \and Feature importance \and
LSTM \and Transformer \and Random Forest}
\end{abstract}

% =======================================================
\section{Introduction}
\label{sec:introduction}

Financial time series prediction is a central task in risk
management, where decisions based on incorrect forecasts can have
quantifiable consequences on portfolios and institutions
\cite{arsenault2025}. In recent years, deep learning models have
reduced prediction error on this type of data compared to classical
statistical methods \cite{bhandari2022,abdulquadir2023}. However,
these models produce their outputs without revealing the factors
that determine them, which restricts their adoption in regulated
environments where transparency is not optional \cite{arsenault2025}.

To address this opacity, the field of Explainable Artificial
Intelligence (XAI) groups methods that identify which factors
influenced a prediction. These methods are organized by mechanism:
gradient-based methods require access to the model's internal
structure, attention-based methods depend on the architecture used,
and perturbation-based methods such as SHAP and LIME are
model-agnostic because they generate explanations by observing how
the output changes when parts of the input are modified
\cite{arsenault2025}. However, these perturbation methods were
formulated under the assumption that input variables are
independent. When applied to time series, they treat each time step
as an isolated variable, ignoring that observations form a sequence
with structural dependencies \cite{tonekaboni2020,huang2025}. In
the financial domain, where market dynamics depend on patterns
accumulated over time, this omission can produce explanations that
do not reflect the actual behavior of the model
\cite{yadav2025,crabbe2021}.

Attention mechanisms in architectures such as LSTM with attention or
Transformers produce weights indicating where the model focused, but
being integrated within the architecture they are only valid for
that specific model and cannot transfer to models of a different
nature. This leaves perturbation-based methods as the only
model-agnostic family, but their application to time series requires
explicitly incorporating the temporal dependency their standard
formulation ignores. Methods that do respect temporal ordering while
remaining model-agnostic, such as dynamic masks \cite{crabbe2021} and
perturbation-window frameworks \cite{roshinta2026}, have not been
evaluated on models of fundamentally different computational natures
trained on the same data, which prevents determining whether their
explanations reflect model behavior or are a consequence of the
chosen architecture.

This paper proposes a model-agnostic temporal explainability
framework that identifies which variables and at which historical
moment were most relevant for a prediction, without depending on
the internal structure of the model used. The framework operates
through temporally ordered perturbations on the input window,
producing an importance matrix per variable and time step that can
be computed over any prediction model. Its model-agnostic nature
is demonstrated by training models of fundamentally different
computational natures on financial time series and verifying that
the generated explanations capture the specific behavior of each
model, rather than being independent of the explained model as in
the failure mode documented by Adebayo et al.~\cite{adebayo2018}.

% =======================================================
\section{Related Work}
\label{sec:related_work}

Related work is organized into two thematic axes. (1) XAI methods
incorporating the temporal dimension. (2) Temporally-aware
model-agnostic explainability frameworks.

\subsection{XAI Methods Incorporating the Temporal Dimension}

Huang et al.~\cite{huang2025} proposed ShapeX, which segments series
into shapelets and computes Shapley values over them for
classification, where a recognizable pattern over a subsequence is
itself the evidence for the predicted label. This work targets
continuous prediction over a multivariate window instead, where that
premise breaks down. Two variables at the same instant can have
independent causal roles, and grouping them into a shared segment
would dilute that distinction, motivating cell-level rather than
segment-level attribution.

Yadav and Subbian~\cite{yadav2025} showed that gradient, occlusion,
and permutation-based methods fail on dynamic prediction because they
do not adapt to time-varying variable-target relationships, producing
attributions that jump abruptly instead of varying smoothly. Their
proposed fix, learnable masks optimized via gradients, ties the
solution to differentiable models. This work shares the diagnosis but
resolves it without that cost, respecting causal order per
perturbation instead of optimizing a global gradient-based mask.

This pattern is not exclusive to generic domains. It repeats in
financial applications combining SHAP or LIME with LSTM/CNN
architectures for auditing, non-technical decision support, or trend
classification \cite{gupta2025,sathyampeth2025,manikrao2026,
jagannathan2025,kumar2024,singh2025,muhammad2024}, applying these
methods directly on the series without adapting their perturbation
mechanism to temporal order.

\subsection{Temporally-Aware Model-Agnostic Explainability Frameworks}

Tonekaboni et al.~\cite{tonekaboni2020} introduced FIT, which
quantifies observation importance via the prediction change upon
removal, controlling for the natural prediction drift as more data
becomes available. Its mechanism needs only model predictions, so it
could in principle run on any model, but it was evaluated on a
single model type, leaving open whether its attributions reflect
model behavior or architecture.

Bento et al.~\cite{bento2021} proposed TimeSHAP, extending KernelSHAP
to sequences with feature, timestep, and cell-level attributions,
evaluated on bank fraud detection, the domain closest to this work.
Its coalition-sampling mechanism over reference sequences restricts
it to natively sequential models, precluding direct use on Random
Forest, and Shapley-sampling cost scales unfavorably with $T$ and
$V$.

Raykar et al.~\cite{raykar2023} proposed TsSHAP, applying TreeSHAP
over a surrogate trained on user-defined interpretable features for
univariate forecasting, without attributing importance to individual
variables within a multivariate window.

Crabbé and Van Der Schaar~\cite{crabbe2021} proposed DynaMask, which
fits a gradient-optimized perturbation mask penalizing abrupt jumps
to guarantee temporal coherence, producing a $T{\times}V$ score and
making it the most direct antecedent of this work. However, it
requires model gradients, restricting it to neural networks.

Leung et al.~\cite{leung2023winit} proposed WinIT, extending
FIT~\cite{tonekaboni2020} with importance per \emph{window} of recent
observations rather than per individual observation, mitigating
diluted attribution among temporally correlated variables. Like FIT,
its mechanism needs no gradients or internal access, but it was
evaluated only under recurrent models, leaving cross-architecture
consistency unverified.

Roshinta et al.~\cite{roshinta2026} proposed a dynamic-window
perturbation framework using PCA to reduce redundancy among
correlated variables before perturbing, reconstructing attributions
in the original variable space. It is mechanically closest to this
proposal, needing no Shapley values nor internal access, but
per-variable attribution is indirect, reconstructed from principal
components rather than obtained per cell, and evaluation is limited
to recurrent architectures (GRU, LSTM).

Kim and Park~\cite{gsshap2026} proposed GroupSegment-SHAP (GS-SHAP),
the work closest in domain to this proposal. It groups correlated
variables via the HSIC criterion, detects change points within each
group via recursive MMD divergence producing segments
$\mathcal{S}^{(k)}$, and computes Shapley values over structured
players $p_{k,j} = G_k \times S_j^{(k)}$ via permutation sampling
with mean-replacement masking, improving structural coherence by
exploiting inter-variable correlation. However, GS-SHAP (1)~is
evaluated on a single model per task, and (2)~its output granularity
operates at the \emph{group} level ($K$ groups $\times$ $J$
segments), losing the per-variable resolution that this framework
produces directly in $\hat{\alpha} \in \mathbb{R}^{T \times V}$ with
$V$ original ungrouped variables.

The works in this axis share a common limitation. None evaluates its
explanations on models of different computational natures trained on
the same data. TimeSHAP and DynaMask cannot be applied to ensemble
models such as Random Forest, the former due to its dependence on
sequential model structure and the latter due to requiring gradients.
TsSHAP operates on predefined features and does not attribute
importance per individual cell. FIT, WinIT, Roshinta et al., and
GS-SHAP evaluate on a single model type, which prevents ruling out
the failure mode documented by Adebayo et al.~\cite{adebayo2018},
attributions that reflect the perturbation method rather than the
explained model's behavior. GS-SHAP aggregates variables into groups
and Roshinta et al. reconstructs them indirectly via PCA, both losing
the direct per-variable resolution that this framework produces. This
work proposes a model-agnostic temporal perturbation framework that
produces importances at full $(T \times V)$ resolution without prior
grouping, can operate on any model through an injectable prediction
function, and whose agnostic nature is demonstrated by evaluating
models of different computational natures on the same data.
Table~\ref{tab:related_work} summarizes the comparison.

\begin{table}[ht]
\caption{Comparison of related work along key problem dimensions.
\cmark~=~yes; \xmark~=~no; $\sim$~=~partially.
\textit{Temp. resp.}~= explanations respect temporal ordering;
\textit{Agnostic}~= no access to model internals required;
\textit{Temp. imp.}~= importance attributed per time step;
\textit{Feat. imp.}~= importance attributed per input feature;
\textit{Agn. verif.}~= behavior evaluated across models of
different computational natures.}
\label{tab:related_work}
\scriptsize
\setlength{\tabcolsep}{3pt}
\resizebox{\textwidth}{!}{%
\begin{tabular}{lccccccl}
\toprule
\textbf{Work} &
\textbf{XAI} &
\textbf{Temp. resp.} &
\textbf{Agnostic} &
\textbf{Temp. imp.} &
\textbf{Feat. imp.} &
\textbf{Agn. verif.} &
\textbf{Type} \\
\midrule
\multicolumn{8}{l}{\textit{Axis 1 --- XAI with temporal dimension}} \\
\midrule
Huang et al.~\cite{huang2025}
& \cmark & \cmark & \cmark & \cmark & \cmark & \xmark & ShapeX \\
Yadav \& Subbian~\cite{yadav2025}
& \cmark & \cmark & \cmark & \cmark & \cmark & \xmark & Mask Eval. \\
\midrule
\multicolumn{8}{l}{\textit{Axis 2 --- Temporally-aware agnostic explainability}} \\
\midrule
Tonekaboni et al.~\cite{tonekaboni2020}
& \cmark & \cmark & \cmark & \cmark & \cmark & \xmark & FIT \\
Leung et al.~\cite{leung2023winit}
& \cmark & \cmark & \cmark & \cmark & \cmark & \xmark & WinIT \\
Crabbé \& Van Der Schaar~\cite{crabbe2021}
& \cmark & \cmark & $\sim$ & \cmark & \cmark & \xmark & DynaMask \\
Bento et al.~\cite{bento2021}
& \cmark & \cmark & $\sim$ & \cmark & \cmark & \xmark & TimeSHAP \\
Raykar et al.~\cite{raykar2023}
& \cmark & \cmark & \cmark & \xmark & $\sim$ & \xmark & TsSHAP \\
Roshinta et al.~\cite{roshinta2026}
& \cmark & \cmark & \cmark & \cmark & $\sim$ & \xmark & Dyn. windows \\
Kim and Park~\cite{gsshap2026}
& \cmark & \cmark & \cmark & \cmark & $\sim$ & \xmark & GS-SHAP \\
\midrule
\textbf{This work}
& \cmark & \cmark & \cmark & \cmark & \cmark & \cmark & Agnostic Fwk. \\
\bottomrule
\end{tabular}}
\end{table}

% =======================================================
\section{Theoretical Background}
\label{sec:background}

\subsubsection{Post-hoc Explainability.}
A method is post-hoc when it operates on an already-trained model
without modifying its structure or training process \cite{arsenault2025}.
Given a trained model and a specific prediction, the post-hoc method
produces an importance attribution over the input variables.

\subsubsection{Model-Agnostic Methods.}
A post-hoc method is model-agnostic when it requires no access to
the model's internal structure. It can be applied to neural networks,
ensemble models, or statistical models without knowing their weights,
gradients, or architecture \cite{arsenault2025}. This property
enables comparing explanations across models under identical
conditions.

Being agnostic by design does not guarantee that a method's
attributions reflect the actual behavior of the explained model.
Adebayo et al.~\cite{adebayo2018} showed that several saliency
methods produce nearly identical explanations even when a model's
weights are randomized, that is, attributions independent of both
the model and the data used to train it. This implies that the
mechanical agnosticism of a method must be verified empirically, and
cannot be assumed from its design alone.

\subsubsection{Perturbation-Based Methods.}
Perturbation methods estimate the importance of a variable by
observing how the model's prediction changes when that variable is
modified or removed from the input \cite{arsenault2025}. SHAP
formalizes this with Shapley values from cooperative game theory;
LIME fits a local linear model over perturbations near the instance.
Both were formulated assuming independent input variables, which
makes them generate temporally inconsistent sequences when applied
to time series \cite{tonekaboni2020,crabbe2021}.

\subsubsection{Temporal Explainability.}
Temporal explainability extends importance attribution to the time
dimension, recognizing that the same variable may be decisive at one
historical moment and irrelevant at another \cite{tonekaboni2020}.
Given a model and an input window of $T$ time steps and $V$
variables, temporal explainability produces a matrix
$\hat{\alpha} \in \mathbb{R}^{T \times V}$, where each cell
$\hat{\alpha}(t, v)$ represents the relevance of variable $v$ at
time step $t$ for the analyzed prediction. For a temporal
explanation to be valid, perturbations must respect temporal
ordering: future information cannot be modified to explain a past
instant \cite{tonekaboni2020,crabbe2021}.

% =======================================================
\section{Proposed Framework}
\label{sec:framework}

The proposed framework produces a matrix
$\hat{\alpha} \in \mathbb{R}^{T \times V}$ where each cell
indicates how much variable $v$ at time step $t$ contributed to
the analyzed prediction. It is composed of four modules with
well-defined responsibilities, illustrated in
Fig.~\ref{fig:pipeline}: (1)~a data ingestion and preparation
module, (2)~a uniform prediction interface that decouples the
perturbation engine from any concrete model, (3)~the temporal
perturbation engine that produces $\hat{\alpha}$, and (4)~a
statistical validation module. Modules (2) and (3) form the
agnostic core: module (2) ensures any model can be connected
without modifying the engine, and module (3) ensures perturbations
are temporally coherent regardless of the connected model.

\begin{figure}[ht]
\centering
\includegraphics[width=\textwidth]{../PIPELINE.png}
\caption{Framework pipeline. Five modules.
(1)~Data ingestion and preparation, S\&P~500 acquisition, $V{=}15$
feature derivation, sliding windows of shape $(N, 20, 15)$.
(2)~Training of three models of different computational natures
(LSTM, Transformer, Random Forest) under the same interface.
(3)~Uniform prediction interface via a two-level adapter pattern
that decouples the engine from any concrete model.
(4)~Temporal perturbation engine producing
$\hat{\alpha} \in \mathbb{R}^{20 \times 15}$ per instance via a
single batch of 301 perturbed windows.
(5)~Statistical validation through inter-model consistency
(Spearman~$\rho$), significance (Mann-Whitney~$U$), and additional
robustness checks (alternative reference, market regime, bounded
target).}
\label{fig:pipeline}
\end{figure}

\subsection{Data Ingestion and Preparation}

The framework receives as input a multivariate time series
organized in fixed-length sliding windows. Historical data of the
S\&P~500 index (\texttt{\^{}GSPC}) for the period 2010--2024 was
obtained via the \textit{yfinance} library, totaling 3,754 trading
days. From five base price and volume variables (Open, High, Low,
Close, Volume), ten technical indicators are derived: simple moving
averages of 5, 10, and 20 days (SMA\_5, SMA\_10, SMA\_20), MACD
and its signal (MACD, MACD\_Signal), 14-period relative strength
index (RSI\_14), Bollinger bandwidth (BB\_width), 14-period average
true range (ATR\_14), log return (Log\_Return), and relative volume
(Volume\_ratio), forming a space of $V=15$ variables.

Data were partitioned respecting temporal order into a development
period (80\%) and test set (20\%). The MinMax scaler was fitted on
the complete development period to avoid data leakage. Sliding
windows of $T=20$ days with prediction horizon $H=1$ day were
constructed, where the target variable is the normalized closing
price of the next day. Each window has shape $(T, V) = (20, 15)$.

\subsection{Uniform Prediction Interface}

Model-agnosticism is implemented through a two-level adapter pattern.
The \texttt{TemporalPerturbationEngine} receives as its sole contact
point with the model a prediction function
$f: \mathbb{R}^{N \times T \times V} \rightarrow \mathbb{R}^N$
injected into its constructor; the engine never accesses the model
directly. This separation ensures that incorporating a new model
requires no modification to any existing component.

\textbf{Level 1, signature adaptation.} Each model implements a
wrapper standardizing the input signature to $(N, T, V)$ and the
output to $(N,)$: the Random Forest wrapper flattens each window
$(T,V) \rightarrow (T{\cdot}V)$ before invoking the scikit-learn
regressor, while LSTM and Transformer wrappers convert NumPy arrays
to PyTorch tensors and run gradient-free inference. From the
engine's perspective, any model behaves identically.

\textbf{Level 2, prediction function injection.} The
\texttt{ModelAgnosticExplainer} builds the callable $f$ from the
wrapper and injects it into the engine, automated for PyTorch and
scikit-learn backends. For any other model, statistical, proprietary,
or from a different framework, the user can inject a function with
signature $(N, T, V) \rightarrow (N,)$ directly, without modifying
any framework component, implementing dependency inversion.

\subsection{Temporal Perturbation Engine}

Given an input window $x \in \mathbb{R}^{T \times V}$ and a model
$f$, the engine constructs an importance matrix
$\alpha \in \mathbb{R}^{T \times V}$ where each cell $\alpha(t, v)$
quantifies how much the model's prediction changes when variable $v$
at time step $t$ is replaced by a neutral reference value:

\begin{equation}
    \alpha(t, v) = \left| f(x) - f(x_{-(t,v)}) \right|
\end{equation}

where $x_{-(t,v)}$ denotes the original window with position
$(t, v)$ replaced by $r_v$, the reference value of variable $v$
defined as its mean over the training set in normalized scale. The
reference vector $\mathbf{r} \in \mathbb{R}^V$ is computed once
before evaluation and kept fixed throughout.

Perturbations respect the temporal ordering: each position $(t, v)$
is masked independently without modifying any other value,
guaranteeing that perturbed sequences are temporally coherent and
do not introduce future information into past instants
\cite{tonekaboni2020,crabbe2021}.

For computational efficiency, the engine builds a single batch of
$T \times V + 1$ windows---the original window plus one perturbed
window per position---obtaining all predictions in a single
inference call:

\begin{equation}
    \text{batch} = \left[ x,\; x_{-(0,0)},\; \ldots,\;
    x_{-(T-1,V-1)} \right] \in
    \mathbb{R}^{(TV+1) \times T \times V}
\end{equation}

With $T=20$ and $V=15$, the batch contains 301 windows per
explained prediction. Importance values are then normalized by the
$L_1$ norm to be comparable across predictions and models:

\begin{equation}
    \hat{\alpha}(t, v) = \frac{\alpha(t, v)}
    {\sum_{t'}\sum_{v'} \alpha(t', v')}
\end{equation}

\subsection{Statistical Validation Module}

\textbf{Inter-model consistency.}
This verifies that the framework produces explanations reflecting
each model's specific behavior and not an artifact of the
perturbation method, in the sense flagged by Adebayo et
al.~\cite{adebayo2018} for model-independent saliency methods. For
each pair of models, the Spearman correlation $\rho$ between the
flattened vectors of their $\hat{\alpha}$ matrices is computed,
averaged over $N=100$ randomly selected test instances:

\begin{equation}
    \rho_{\text{inter}} = \frac{1}{N} \sum_{i=1}^{N}
    \rho\!\left(
    \text{vec}(\hat{\alpha}_i^{(A)}),\;
    \text{vec}(\hat{\alpha}_i^{(B)})
    \right)
\end{equation}

High correlation between models of the same family (neural) and
lower correlation between different families (neural vs.\ ensemble)
indicates that explanations capture model-specific behavior.

\textbf{Statistical significance.}
Explanations are compared against a random baseline: matrices of
shape $(T, V)$ with values drawn from a uniform distribution and
$L_1$-normalized, maintaining the same scale as real $\hat{\alpha}$
matrices. The non-parametric Mann-Whitney U test is applied between
real and baseline distributions without assuming normality
\cite{schlegel2019}.

% =======================================================
\section{Experimental Results}
\label{sec:results}

\subsection{Experimental Setup}

S\&P~500 historical data (\texttt{\^{}GSPC}) for 2010--2024
(3,754 trading days, $V{=}15$ variables: 5 base price/volume series
plus 10 derived technical indicators) was partitioned respecting
temporal order into development (80\%) and test (20\%) sets.
Sliding windows of $T{=}20$ days with one-day-ahead prediction
horizon were constructed (shape $(T,V){=}(20,15)$); a MinMax scaler
fitted on the development set was applied to prevent data leakage.
The reference vector $\mathbf{r}{\in}\mathbb{R}^{15}$ is the
per-variable mean over the training set in normalised scale,
fixed for all evaluations.

Three models of fundamentally different computational natures were
trained on the same data under identical conditions:

\begin{itemize}\setlength\itemsep{0pt}
  \item \textbf{LSTM:} 2-layer recurrent network
    (hidden$=$128, dropout~0.2).
  \item \textbf{Transformer:} 2-layer encoder with 4 attention heads,
    feedforward dimension 256, sinusoidal positional encoding,
    and learned attention pooling over the full sequence.
  \item \textbf{Random Forest:} 200 decision trees
    (\texttt{RandomForestRegressor}), non-sequential tabular baseline.
\end{itemize}

All neural models were trained with AdamW ($lr{=}10^{-3}$),
learning rate reduction on plateau (factor~0.5, patience~5),
gradient clipping (norm~1.0), and early stopping (patience~15).
Test-set MSE: LSTM~$0.00342$, Transformer~$0.00284$, RF~$0.12131$
(higher RF error attributable to price extrapolation beyond the
training-period maximum during the 2023--2024 bull market).

As an additional reference, a naive persistence baseline
$f(x)=x_{t,\text{Close}}$ (predicting no change) achieves
MSE$=0.00016$, an order of magnitude below all three trained models.
This does not invalidate the cross-model comparison, since this
work's goal is verifying framework agnosticism rather than optimizing
predictive accuracy, but it indicates none of the three models beats
the random-walk hypothesis on nominal price prediction. It also
connects to a scale artefact. The MinMax scaler is fit on the
development period (2010--2021), but the S\&P~500 reached new
all-time highs through most of the test period (2022--2024, 731
windows), so $32.7\%$ of real test values fall outside $[0,1]$.
Random Forest cannot extrapolate beyond its training-observed output
range by design ($0.00\%$ of its predictions fall outside $[0,1]$,
collapsing to $[0.645,0.657]$), inflating its MSE beyond the
architectural gap alone. Section~\ref{sec:robustness} evaluates a
naturally bounded target that resolves this.

\subsection{Advance 1: Cell-Level Attribution Without Gradient Access}
\label{sec:advance1}

\textbf{Gap from DynaMask.}
DynaMask~\cite{crabbe2021} is the closest antecedent in output
format: it produces a $T{\times}V$ importance matrix via a
gradient-optimised perturbation mask, reporting AUP\,$=0.97\pm0.01$
on a synthetic benchmark (AR(1) inputs, $f{=}\Sigma x^2$,
$T{=}50$, $V{=}4$; their Table~1).
The fundamental restriction is that gradient access to the target
model is mandatory; DynaMask cannot be applied to gradient-free
models such as Random Forest.

\textbf{This work.}
FO computes attributions via a single batched call of
$T{\times}V{+}1{=}301$ perturbed windows, requiring only
prediction-function access. We replicate the benchmark family
($T{=}20$, $V{=}15$, $n{=}10$ repetitions) under two ground-truth
scenarios:

\begin{itemize}\setlength\itemsep{0pt}
  \item \textit{Rare Feature (RF):} $|\mathcal{A}_V|{=}3$
    salient variables, all $T$ steps relevant.
  \item \textit{Rare Time (RT):} $|\mathcal{A}_T|{=}4$ salient
    time steps, all $V$ variables relevant.
\end{itemize}

The white-box target is $f(\mathbf{X}) = \sum_{(t,v)\in\mathcal{A}}
x_{t,v}^{2}$; inputs are independent AR(1) per variable.
Metrics follow~\cite{crabbe2021}: AUP ($\uparrow$), AUR ($\uparrow$),
mask information content $I_M$ ($\uparrow$), mask entropy $S_M$
($\downarrow$).

\begin{table}[ht]
\caption{Synthetic white-box validation ($n{=}10$ repetitions,
         $N{=}30$; SVS: $N{=}5$\,\textsuperscript{$\star$}).}
\label{tab:exp1_synthetic}
\centering
\scriptsize
\setlength{\tabcolsep}{4pt}
\begin{tabular}{lcccccccc}
\toprule
& \multicolumn{4}{c}{\textbf{Rare Feature (RF)}}
& \multicolumn{4}{c}{\textbf{Rare Time (RT)}} \\
\cmidrule(lr){2-5}\cmidrule(lr){6-9}
\textbf{Method} & \textbf{AUP} & \textbf{AUR}
  & $I_M$ & $S_M$
  & \textbf{AUP} & \textbf{AUR}
  & $I_M$ & $S_M$ \\
\midrule
FO (proposed)
  & $1.00{\pm}0.00$ & $0.16{\pm}0.01$ & $30.6{\pm}1.1$ & $0.00{\pm}0.00$
  & $1.00{\pm}0.00$ & $0.16{\pm}0.01$ & $30.0{\pm}0.8$ & $0.00{\pm}0.00$ \\
FP
  & $1.00{\pm}0.00$ & $0.18{\pm}0.01$ & $31.9{\pm}1.2$ & $0.00{\pm}0.00$
  & $1.00{\pm}0.00$ & $0.18{\pm}0.01$ & $32.1{\pm}1.6$ & $0.00{\pm}0.00$ \\
IG
  & $1.00{\pm}0.00$ & $0.16{\pm}0.01$ & $30.6{\pm}1.1$ & $0.00{\pm}0.00$
  & $1.00{\pm}0.00$ & $0.16{\pm}0.01$ & $30.0{\pm}0.8$ & $0.00{\pm}0.00$ \\
SVS\textsuperscript{$\star$}
  & $1.00{\pm}0.00$ & $0.17{\pm}0.01$ & $31.2{\pm}0.8$ & $0.00{\pm}0.00$
  & $1.00{\pm}0.00$ & $0.15{\pm}0.01$ & $28.9{\pm}1.5$ & $0.00{\pm}0.00$ \\
\bottomrule
\end{tabular}
\end{table}

FO achieves AUP\,$=0.9960\pm0.0000$ in both scenarios, exceeding the
DynaMask FO baseline of $0.97$ while requiring no gradient access and
running on LSTM, Transformer, \emph{and} Random Forest without
modification. $S_M \approx 4.4{\times}10^{-5}$ is a floating-point
residual of $L_1$ normalisation (analytically zero, since cells
outside $\mathcal{A}$ receive zero raw importance). This confirms
per-cell resolution over all $V{=}15$ original ungrouped variables,
the $T{\times}V$ output format that prior gradient-free methods had
not demonstrated simultaneously across neural and ensemble backends.

SVS (Shapley Value Sampling) is the only other method that is
model-agnostic, produces a $T{\times}V$ matrix, and requires no
gradient access --- making it the valid same-conditions baseline.
Both FO and SVS achieve AUP\,$=0.9960\pm0.0000$ in both scenarios,
confirming Shapley-equivalent identification precision on a
ground-truth benchmark. AUR differences remain within one standard
deviation. The substantive difference between the two methods is
therefore computational: quantified in Sections~\ref{sec:advance2}
and~\ref{sec:results_faithfulness}.

\subsection{Advance 2: Order-of-Magnitude Runtime Reduction}
\label{sec:advance2}

\textbf{Gap from TimeSHAP and GS-SHAP.}
Two prior works share the S\&P~500, $T{=}20$ domain and report
per-sample runtimes under a controlled perturbation budget of
$M{=}50$ permutations (Kim and Park~\cite{gsshap2026}, Table~C3,
CPU): TimeSHAP~\cite{bento2021} at \textbf{2,813\,ms/sample},
requiring a reference sequence and sequential model structure that
precludes Random Forest, and GS-SHAP~\cite{gsshap2026} at
\textbf{1,194\,ms/sample}, model-agnostic but outputting $K{\times}J$
group-segment players evaluated on a single BiLSTM. Both runtimes
scale with the number of coalition evaluations required by Shapley
value computation.

\textbf{This work.}
FO reduces the inference count to a single batch of
$T{\times}V{+}1{=}301$ windows per instance, no coalition sampling,
no permutations, deterministic result.
Under the same domain (S\&P~500, $T{=}20$, CPU), FO achieves
\textbf{17.6\,ms/sample} (LSTM), a reduction of ${\times}159.8$
over TimeSHAP and ${\times}67.8$ over GS-SHAP.
Table~\ref{tab:advances} summarises all three concrete advances over
prior work with verifiable numbers.

\begin{table}[ht]
\caption{Concrete advances over prior work.
  $\dagger$: reported by Kim and Park~\cite{gsshap2026}, Table~C3
  (S\&P~500, $T{=}20$, $M{=}50$, CPU).
  $\ddagger$: measured in this work, same domain and $T$.
  $\star$: different synthetic setup ($T{=}50$, $V{=}4$) vs.\
  ours ($T{=}20$, $V{=}15$); same function class.}
\label{tab:advances}
\centering
\scriptsize
\setlength{\tabcolsep}{4pt}
\begin{tabular}{lp{2.8cm}cc}
\toprule
\textbf{Prior work} & \textbf{Open gap} &
\textbf{Their result} & \textbf{This work} \\
\midrule
DynaMask~\cite{crabbe2021}
  & Requires gradients; no RF support
  & AUP $0.97^{\star}$, no RF
  & AUP $\mathbf{0.9960}$, RF \cmark \\[2pt]
TimeSHAP~\cite{bento2021}
  & Sequential models only; slow
  & $2{,}813$\,ms/sample$^{\dagger}$
  & $\mathbf{17.6}$\,ms/sample$^{\ddagger}$ (${\times}159.8$) \\[2pt]
GS-SHAP~\cite{gsshap2026}
  & Grouped output ($K{\times}J$); single model
  & $1{,}194$\,ms/sample$^{\dagger}$; $K{\times}J$ only
  & $\mathbf{17.6}$\,ms/sample$^{\ddagger}$; $T{\times}V$ \cmark \\
\bottomrule
\end{tabular}
\end{table}

\subsection{Model-Agnosticism Verification}
\label{sec:agnosticism}

A key claim of this work is that the framework generates
\emph{model-specific} explanations --- attributions that reflect the
internal behavior of the queried model, not an artefact of the
perturbation method. This is evaluated through two complementary
tests on 100 randomly selected test instances.

\textbf{Inter-model consistency (Spearman~$\rho$).}
For each pair of models, we compute the mean Spearman $\rho$ between
their flattened $\hat{\alpha}$ vectors over all instances
(Table~\ref{tab:consistency}).

\begin{table}[ht]
\caption{Inter-model consistency (Spearman $\rho$, mean over 100
instances). High within-family, lower across-family correlation
confirms that attributions capture model-specific behavior.}
\label{tab:consistency}
\centering
\begin{tabular}{lc}
\toprule
\textbf{Model pair} & \textbf{Spearman $\rho$} \\
\midrule
LSTM -- Transformer          & 0.748 \\
LSTM -- Random Forest        & 0.255 \\
Transformer -- Random Forest & 0.382 \\
\bottomrule
\end{tabular}
\end{table}

The LSTM--Transformer pair yields $\rho{=}0.748$: both neural
architectures agree on the most salient importance patterns despite
differing mechanisms. Neural--RF pairs yield $\rho{<}0.40$, reflecting
qualitatively different information processing by the ensemble.
This cross-family divergence is only observable under a framework
that runs on all three architectures simultaneously; DynaMask requires
gradients (precluding RF) and TimeSHAP requires sequential structure
(precluding RF).

\textbf{Statistical significance (Mann-Whitney $U$).}
Attributions are compared against $L_1$-normalised uniform random
matrices of identical shape using the Mann-Whitney $U$ test. All
three models reject the null hypothesis of equivalence with the
random baseline by tens of orders of magnitude ($U{=}3.55{\times}10^9$
LSTM, $3.59{\times}10^9$ Transformer, $2.72{\times}10^8$ Random
Forest; $p{<}10^{-300}$ in all three cases), confirming that the
framework produces informative, non-trivial attributions regardless
of model nature.

\subsection{Robustness Checks}
\label{sec:robustness}

The tests above use a single reference vector, a single target
formulation, and a random test sample. This subsection reports three
further checks reusing the checkpoints already described, plus a
log-return retraining.

\textbf{Naturally bounded target.} Retraining all three models on
next-day log return instead of \texttt{Close} (same split, scaler,
architectures) removes the MinMax saturation flagged in
Section~\ref{sec:advance1}. Under this bounded target, \emph{no}
model beats a zero-return persistence baseline (MSE$=0.00257$; LSTM
$0.00268$, Transformer $0.00260$, RF $0.00446$), consistent with
weak-form market efficiency. The RF-vs-neural MSE gap shrinks from
${\sim}35{\times}$ to ${\sim}1.7{\times}$, and inter-model consistency
reproduces the same pattern (LSTM--Transformer $\rho{=}0.541$,
LSTM--RF $\rho{=}0.187$, Transformer--RF $\rho{=}0.311$; $p{\approx}0$
vs.\ random), confirming agnosticism verification is not an artefact
of price autocorrelation.

\textbf{Directional accuracy.} Direction-of-movement accuracy against
the previous close, for both targets, falls within a narrow band
around chance for all six model-target combinations (Close: LSTM
$47.1\%$, Transformer $46.4\%$, RF $48.0\%$; Log\_Return: LSTM
$46.6\%$, Transformer $50.3\%$, RF $47.5\%$), reinforcing the
market-efficiency reading above.

\textbf{Reference sensitivity and market regime.} Table~\ref{tab:robustness}
reports inter-model consistency under two alternative references
(median, zero) on the same 100 instances, and under the two market
regimes present in the test period (2022 bear, 188 windows;
2023--2024 bull, 501 windows; the test period excludes the 2020
COVID crash, which falls in the development split).

\begin{table}[ht]
\caption{Inter-model consistency (Spearman $\rho$) across reference
vectors and market regimes. The neural--neural pair exceeds both
neural--ensemble pairs in every row.}
\label{tab:robustness}
\centering
\begin{tabular}{lccc}
\toprule
\textbf{Condition} & \textbf{LSTM--Transf.} & \textbf{LSTM--RF} & \textbf{Transf.--RF} \\
\midrule
Mean reference (baseline)   & 0.748 & 0.255 & 0.382 \\
Median reference            & 0.750 & 0.332 & 0.475 \\
Zero reference               & 0.587 & 0.255 & 0.370 \\
\midrule
Bear market 2022             & 0.703 & 0.173 & 0.287 \\
Bull market 2023--2024       & 0.764 & 0.278 & 0.404 \\
\bottomrule
\end{tabular}
\end{table}

The neural-vs-ensemble gap holds under every reference and every
regime, confirming that the framework's central finding does not
depend on the arbitrary choice of reference vector nor on the market
window sampled. Consistency is systematically lower under the bear
regime for all three model pairs, suggesting the three architectures
agree less on what drives predictions when the market is more
volatile. LSTM and Transformer attributions are highly stable across
reference choices (mean-vs-median $\rho{\approx}0.99$); Random
Forest is markedly more sensitive (mean-vs-median $\rho{=}0.81$,
mean-vs-zero $\rho{=}0.67$).

\subsection{Same-Conditions Quality Benchmark: FO vs.\ SVS}
\label{sec:results_faithfulness}

To quantify attribution quality in the absence of ground truth, we
apply the faithfulness protocol of~\cite{crabbe2021}: for each
fraction $\alpha{\in}\{0.1,\ldots,0.6\}$, the
$\lfloor\alpha{\cdot}T{\cdot}V\rfloor$ highest-attribution cells are
replaced by their reference value and the mean absolute prediction
shift is measured:

\begin{equation}
    \Delta_\alpha = \frac{1}{N}\sum_{i=1}^{N}
        \left|f(\mathbf{X}_i) - f(\tilde{\mathbf{X}}_i^{(\alpha)})\right|
    \label{eq:faithfulness}
\end{equation}

A faithful explanation produces high $\Delta_\alpha$: perturbing
truly relevant cells significantly affects model predictions.

SVS (Shapley Value Sampling, 15~permutations) is the only other
method that is model-agnostic, produces a $T{\times}V$ matrix, and
runs on all three models, making it the valid same-conditions
comparator. Evaluation uses the same S\&P~500 test set ($N{=}100$;
SVS over $N{=}10$ due to computational cost), same reference vector,
same three models.

\begin{table}[ht]
\caption{Prediction shift $\Delta_{0.3}$ at $\alpha{=}0.3$.
         IG not applicable to Random Forest ($\dagger$);
         SVS: $N{=}10$, 15~permutations. Bold: best per model.}
\label{tab:faithfulness}
\centering
\begin{tabular}{lcccc}
\toprule
\textbf{Model} & \textbf{FO} & \textbf{FP}
               & \textbf{IG} & \textbf{SVS} \\
\midrule
LSTM          & \textbf{0.5033} & 0.4776 & \textbf{0.5034} & 0.4891 \\
Transformer   & 0.4518          & 0.4320 & \textbf{0.4565} & 0.4352 \\
Random Forest & 0.3047          & 0.2723 & $\dagger$        & \textbf{0.3212} \\
\midrule
\multicolumn{5}{l}{\small FO$-$SVS: LSTM $+0.0142$ (FO better)\quad
  Transformer $+0.0166$ (FO better)\quad RF $-0.0165$ (SVS better)} \\
\bottomrule
\end{tabular}
\end{table}

\begin{table}[ht]
\caption{Per-sample runtime (ms, CPU, mean$\,\pm\,$std over 5~runs of
         50~samples). FO: $T{\times}V{+}1{=}301$ forward passes.
         SVS: $2{\times}15{\times}T{\times}V{=}9{,}000$ evaluations.}
\label{tab:timing}
\centering
\begin{tabular}{lccc}
\toprule
\textbf{Model} & \textbf{FO (ms)} & \textbf{SVS (ms)}
               & \textbf{FO speedup} \\
\midrule
LSTM          & $17.6\pm0.9$   & $469.1\pm3.1$   & ${\times}26.6$ \\
Transformer   & $13.8\pm0.3$   & $413.5\pm15.9$  & ${\times}30.0$ \\
Random Forest & $112.0\pm3.9$  & $2519.4\pm82.4$ & ${\times}22.5$ \\
\bottomrule
\end{tabular}
\end{table}

On neural models, FO surpasses SVS in faithfulness ($\Delta_{0.3}$:
$+0.014$ LSTM, $+0.017$ Transformer) while running $22$--$30{\times}$
faster.
On Random Forest, SVS outperforms FO by $0.017$: coalition sampling
captures combinatorial interactions among decision-tree splits that
single-cell occlusion misses.
The structural source of FO's speed advantage is that it issues
exactly one batched inference call of size $T{\times}V{+}1{=}301$,
whereas SVS issues $2{\times}n_\text{perms}{\times}T{\times}V{=}9{,}000$
evaluations at 15~permutations.
IG surpasses FO on Transformer by $0.005$ at the cost of gradient
access (precluding RF).
FP consistently ranks last: replacing a cell with another observed
value introduces exploitable market signal, artificially attenuating
the measured perturbation impact.

% =======================================================
\section{Conclusions}
\label{sec:conclusions}

This paper presented a model-agnostic temporal explainability
framework for financial time series prediction. Its core contribution
is a two-level adapter pattern that decouples the temporal perturbation
engine from any specific model backend, enabling any model that exposes
a prediction function with signature $(N, T, V) \rightarrow (N,)$ to
be connected without modifying any framework component.

Under identical conditions, same test set, reference vector, and
three models, FO matches SVS, the only other model-agnostic method
producing a $T{\times}V$ matrix on all three models, in synthetic
ground-truth precision (AUP $=0.9960$ both) while running
$22$--$30{\times}$ faster on neural architectures and slightly behind
on Random Forest ($-0.017$ faithfulness), where coalition sampling
better captures combinatorial tree interactions.

The framework also produces statistically non-trivial attributions
($p < 10^{-300}$) and inter-model consistency coherent with
architectural families ($\rho = 0.748$ between neural models vs.\
$\rho < 0.40$ against Random Forest), a pattern that holds under
alternative reference vectors, under both market regimes present in
the test period, and under a naturally bounded target that removes
the MinMax saturation affecting Random Forest under \texttt{Close}
(Section~\ref{sec:robustness}). Neither target formulation beats a
naive persistence baseline, in magnitude or in direction, consistent
with weak-form market efficiency. The main remaining limitation is
evaluation on a single financial index. Future work includes
multi-index evaluation and a deeper analysis of why inter-model
consistency decreases under volatile market conditions.

\begin{credits}
\subsubsection{\discintname}
The authors have no competing interests to declare that are
relevant to the content of this article.
\end{credits}

\bibliographystyle{splncs04}
\bibliography{biblio}

\end{document}
